% =====================================================================
\section{실험결과}

\subsection{결과의 범위}
아래 결과는 학습과 겹치지 않는 난수시드 9,900,950으로 만든 30일 연속 진단 실행의
가운데 28일이다. 가중치를 고정하고 탐색확률을 0으로 설정했으며, 같은 날짜의 모든
정책은 같은 트럭·선박·서비스 난수를 공유했다. 결과는 문헌 기반 합성환경에서 관찰된
아키텍처의 가능성을 나타낸다.

\subsection{아키텍처 작동 확인}
구조적 작동은 여섯 조건으로 확인했다. 확인 경로가 조건마다 다르므로 출처를 함께
적는다. 성능 수치(\S5.3 이하)는 모두 판정 실행에서 나오지만, 아래 일부는 자동 시험과
학습 실행에서 확인된 것이다.

제안과 수락을 서로 다른 모형이 판단하는 것은 구현 구조에서, 한쪽 동의만으로 확정된
재배치가 0건인 것은 확정 절차의 제약식~(11)에서 보장된다. 미래 실제 도착·완료시각이
정책 입력으로 들어간 경우가 0건인 것은 구현 검토와 정책별 정보등급 표시로 확인했다.
고정정책 평가 중 반사실 시뮬레이션 호출이 \textbf{0회}인 것은 \emph{판정 실행}에서
직접 확인했다. 행동을 바꾸지 않은 분기가 본 실행의 행동과 비용을 재현하는 동일성
검사 4/4 통과는 \emph{학습 실행}의 결과다---평가에는 반사실 라벨이 없어 이 검사가
성립하지 않는다. 날짜 경계에서 재고·대기열·미완료 작업이 초기화되지 않는 것은 자동
시험이 확인한다.

이는 제안--수락--확정--실행과 반사실 학습--운영 추론의 분리가 하나의 실행 경로로
작동했음을 보여 준다. 판정 실행은 반사실 호출 수·정책 예외·확정 실패 건수를 결과
파일에 함께 남기므로, 이후 실행에서는 실행 의존 항목이 그 실행의 산출물로 확인된다.

\subsection{정책별 날짜 짝비교}
비용 차이는 제안 정책에서 비교 정책을 뺀 값이다. 음수이면 제안 정책의 비용이 더
낮다. 날짜를 독립 비교 단위로 두고 양측 부호검정을 사용했다(표~\ref{tab:paired}).

\begin{table}[t]
\centering
\caption{28일 날짜 짝비교. 음수는 제안 정책의 비용이 낮음을 뜻한다.}
\label{tab:paired}
\small
\begin{tabular}{@{}lrrr@{}}
\toprule
비교 대상 & 제안 정책이 낮은 날 & $p$값 & 날짜 중앙 비용 차이 \\
\midrule
재배치 없음 & \textbf{26/28} & \textbf{$<$0.001} & $-$533만원 \\
선착순 재배치 & 22/28 & 0.0037 & $-$221만원 \\
순이득 기준 재배치 & 21/28 & 0.0125 & $-$246만원 \\
최단 처리시간 재배치 & 21/28 & 0.0125 & $-$212만원 \\
최소 여유시간 재배치 & 20/28 & 0.0357 & $-$250만원 \\
학습 전 모형 & \textbf{18/28} & \textbf{0.1849} & $-$126만원 \\
\bottomrule
\end{tabular}
\end{table}

재배치 없음의 28일 총비용은 100.0억원, 제안 정책은 85.2억원으로 14.71억원(14.7\%)의
감소 방향이 관찰됐다(그림~\ref{fig:policies}). 위 규칙 네 정책은 블록 재배치만
사용하므로 이 비교는 아키텍처 전체의 가능성을 나타낸다. 행동 범위를 맞춘 비교는
\S5.6에서 따로 다룬다. 날짜별 차이의 분포는 그림~\ref{fig:paired}에 있다.

\begin{figure}[t]
\centering
\includegraphics[width=0.9\textwidth]{fig4-policies.pdf}
\caption{정책별 28일 총비용 감소율. 재배치 없음 대비이며 양수가 비용 감소다.}
\label{fig:policies}
\end{figure}

\begin{figure}[t]
\centering
\includegraphics[width=0.9\textwidth]{fig5-paired.pdf}
\caption{날짜별 비용 차이(제안 정책 $-$ 재배치 없음). 28일을 차이 크기로 정렬했으며
음수가 제안 정책의 비용이 낮은 날이다.}
\label{fig:paired}
\end{figure}

\subsection{수요수준별 비용 방향}
전체 비용 감소의 90.5\%가 수요 12,500대와 15,000대인 나흘에서 발생했다
(표~\ref{tab:byload}). 이는 제안 아키텍처를 상시 개입보다 \emph{혼잡이 관측될 때
선택적으로 개입하는 정책}으로 발전시킬 근거를 제공한다. 혼잡 조건의 반복 수를 늘리면
효과 크기의 통계적 확실성도 함께 높일 수 있다.

\begin{table}[t]
\centering
\caption{수요수준별 비용과 감소 비중}
\label{tab:byload}
\small
\begin{tabular}{@{}rrrrrr@{}}
\toprule
일일 수요 & 일수 & 재배치 없음 & 제안 정책 & 비용 감소 & 전체 감소 중 \\
\midrule
3,500 & 12 & 8.2억 & 7.8억 & 0.47억 & 3.2\% \\
5,000 & 8 & 8.5억 & 7.9억 & 0.60억 & 4.1\% \\
7,500 & 4 & 10.0억 & 9.7억 & 0.32억 & 2.2\% \\
\textbf{12,500} & \textbf{2} & 26.0억 & 19.9억 & \textbf{6.11억} & \textbf{41.5\%} \\
\textbf{15,000} & \textbf{2} & 47.2억 & 40.0억 & \textbf{7.22억} & \textbf{49.0\%} \\
\bottomrule
\end{tabular}
\end{table}

\subsection{블록과 도착시각 행동의 분해}
같은 28일에서 두 행동 중 하나만 허용했다(표~\ref{tab:actions}, 그림~\ref{fig:actions}).
도착시각 조정은 혼잡한 나흘에서 13.20억원을 줄였고 28일 전체 감소의 91.4\%를
설명했다. 블록 재배치만 사용한 조건은 원활·보통 수요에서 1.24억원을 줄였으나 가장
높은 수요에서 비용이 증가해 합계가 1.31억원 높았다. 이 분해는 정책이 수요 수준에
따라 공간과 시간 행동의 사용 비중을 조절하는 방향으로 최적화될 수 있음을 보여 준다.

\begin{table}[t]
\centering
\caption{행동 유형별 비용 감소(억원). 양수는 감소, 음수는 증가를 뜻한다.}
\label{tab:actions}
\small
\begin{tabular}{@{}rrrrrr@{}}
\toprule
일일 수요 & 일수 & 두 행동 모두 & 도착시각만 & 블록만 & 학습 전 모형 \\
\midrule
3,500 & 12 & 0.47 & $-$0.13 & \textbf{0.57} & $-$0.09 \\
5,000 & 8 & 0.60 & 0.29 & \textbf{0.67} & 0.31 \\
7,500 & 4 & 0.32 & 0.08 & 0.00 & 0.81 \\
12,500 & 2 & 6.11 & \textbf{5.73} & $-$0.14 & 4.67 \\
15,000 & 2 & 7.22 & \textbf{7.47} & \textbf{$-$2.41} & 2.27 \\
\midrule
\textbf{합계} & \textbf{28} & \textbf{14.71} & \textbf{13.44} & \textbf{$-$1.31} & \textbf{7.97} \\
\bottomrule
\end{tabular}
\end{table}

\begin{figure}[t]
\centering
\includegraphics[width=0.9\textwidth]{fig6-actions.pdf}
\caption{수요수준별·행동유형별 비용 감소. 재배치 없음 대비이며 같은 28일이다.}
\label{fig:actions}
\end{figure}

시간대별 예약 정원은 이 진단에서 개방형이었다. 따라서 도착시각 조정의 수치는 시간
이동 기능이 제공할 수 있는 \emph{가능성 범위}이며, 유한 정원을 넣은 민감도 분석으로
적용 범위를 정교화할 수 있다.

\subsection{같은 행동 범위를 쓰는 규칙 정책과의 비교}
\S5.5는 비용 감소의 대부분이 도착시각 조정에서 나옴을 보였다. 그렇다면 규칙 정책이
같은 행동을 쓰면 어떤가. 공간 축에서 이미 사용하는 \emph{가장 한산한 곳} 기준을 시간
축으로 옮겨 두 규칙 정책을 추가했다. 새 파라미터를 도입하지 않고 기준의 적용 축만
바꾼 것이며, 다른 규칙 정책과 같은 혼잡 조건을 공유한다.

\begin{table}[t]
\centering
\caption{행동 범위별 28일 비용 감소율과 짝비교}
\label{tab:width}
\small
\begin{tabular}{@{}llr@{}}
\toprule
정책 & 사용 행동 & 비용 감소율 \\
\midrule
제안 정책 & 블록 + 도착시각 & \textbf{+14.72\%} \\
제안 정책 (도착시각만) & 도착시각 & +13.44\% \\
학습 전 모형 & 블록 + 도착시각 & +7.97\% \\
\textbf{규칙: 한산한 시간대} & 도착시각 & \textbf{+3.41\%} \\
\textbf{규칙: 한산한 블록·시간대} & 블록 + 도착시각 & \textbf{+1.51\%} \\
규칙: 최소 여유시간 & 블록 & +0.53\% \\
규칙: 선착순 & 블록 & +0.17\% \\
제안 정책 (블록만) & 블록 & $-$1.31\% \\
\midrule
\multicolumn{3}{@{}l@{}}{\textit{짝비교 (앞 정책이 낮은 날 / $p$값 / 날짜 중앙 차이)}} \\
제안 정책 vs 규칙(블록+시각) & \multicolumn{2}{r@{}}{\textbf{20/28} \quad \textbf{0.0357} \quad $-$166만원} \\
제안 정책(시각만) vs 규칙(시각) & \multicolumn{2}{r@{}}{8/28 \quad 0.0357 \quad $+$106만원} \\
규칙(시각) vs 재배치 없음 & \multicolumn{2}{r@{}}{21/28 \quad 0.0125 \quad $-$235만원} \\
규칙(블록+시각) vs 재배치 없음 & \multicolumn{2}{r@{}}{23/28 \quad 0.0009 \quad $-$424만원} \\
\bottomrule
\end{tabular}
\end{table}

세 가지가 관찰된다. 첫째, 규칙만으로도 도착시각 조정은 비용을 유의하게 낮춘다
($+3.41\%$, $p=0.0125$). 조정이라는 행동 자체가 값을 가진다. 둘째, 행동 범위를 맞춘
비교에서 제안 정책이 규칙 정책보다 낮았다(20/28, $p=0.0357$, 총액 14.72\% 대 1.51\%).
셋째, 도착시각 축만 놓고 보면 \emph{날짜 기준으로는 규칙 정책이 더 자주 낮았다}
(제안 정책 8/28). 그러나 총액은 제안 정책이 약 4배 낮았다.

세 번째 관찰의 이유는 수요수준별 분해에 있다. 원활·보통 수요에서는 규칙 정책이
근소하게 낮고(수요 3,500에서 제안 $-0.13$억 대 규칙 $0.06$억), 혼잡한 나흘에서 제안
정책이 3.8--5.6배 낮다(수요 15,000에서 7.47억 대 1.33억). 학습된 선택은 \emph{자주
이기는 방향이 아니라 비용이 큰 날에 크게 낮추는 방향}으로 최적화됐다. 두 지표를 함께
보고해야 이 성질이 드러난다.

같은 분해에서 규칙(블록+시각)은 혼잡 수요에서 0.38억과 0.09억으로 낮아진다. \S5.5에서
관찰한 \emph{블록 재배치가 혼잡에서 비용을 올리는 성질}이 규칙 정책에서도 같은
방향으로 재현됐다. 이는 특정 정책의 특성이 아니라 행동 자체의 성질임을 시사한다.

\subsection{학습 전후 비교와 작업순서 민감도}
재배치 없음 100.0억원에서 학습 전 모형은 92.0억원, 학습된 모형은 85.2억원이었다.
학습 뒤 추가 감소는 6.7억원으로 전체 감소의 약 45\%에 해당한다. 날짜별로는 학습된
모형이 18/28일 더 낮았고 양측 부호검정은 $p=0.1849$였다. 즉 금액과 승패 모두 낮은
비용 방향을 보였으며, 현재 표본에서 통계적 확실성은 제한적이다. 특히 큰 차이가
혼잡한 소수 날짜에 집중됐으므로, 혼잡 날짜를 더 포함한 반복평가가 학습 기여의 범위를
정교하게 보여 줄 수 있다.

크레인 작업순서 규칙은 수요에 따라 상대 순위가 뒤집혔다(그림~\ref{fig:crane}). 현행
최단 처리시간 규칙은 원활·보통 수요에서 다른 규칙보다 22--25\% 높은 비용을 보였으나
수요 7,500 이상에서는 가장 낮았다. 설계된 수요 비중으로 가중한 30일 기대비용은 현행
규칙이 가장 낮다. 이 결과는 재배치 정책의 효과도 작업순서 규칙별로 다시 검증해야
함을 보여 준다.

\begin{figure}[t]
\centering
\includegraphics[width=0.9\textwidth]{fig7-crane.pdf}
\caption{크레인 작업순서 규칙의 수요 민감도. 재배치 없음 조건이며 각 수요에서 가장
낮은 총비용을 1.0으로 둔 상대값이다.}
\label{fig:crane}
\end{figure}

\subsection{연구질문별 관찰}
RQ1에서는 상호 동의와 용량 확인을 거친 재배치가 연속 환경에서 실행됐다. RQ2에서는
평가 중 반사실 호출 없이 두 비용모형만으로 결정했다. RQ3에서는 재배치 없음 대비
대부분의 날짜에서 낮은 비용 방향이 관찰됐다. RQ4에서는 비용 변화가 도착시각 조정과
혼잡일에 집중됐고, 행동 범위를 맞춘 규칙 정책 대비 낮은 비용(20/28, $p=0.0357$)이
관찰됐으나 학습 전 모형 대비는 방향은 같아도 현재 표본에서 확실성이 제한적이었다.
